\documentclass[11pt]{article}
\usepackage[utf8]{inputenc}
\usepackage[T2A]{fontenc}
\usepackage{amsmath,amssymb}
\usepackage{booktabs}
\usepackage{graphicx}
\usepackage{hyperref}
\usepackage{natbib}
\usepackage{multirow}

\title{Morphology-Aware Grammatical Error Correction for Kazakh: \\
A Dual-Model Pipeline with Language-Specific Error Taxonomy}

\author{Saken Tukenov \\
\texttt{gorgtown@gmail.com}}

\date{}

\begin{document}
\maketitle

\begin{abstract}
We present the first comprehensive GEC system for Kazakh, an agglutinative Turkic language with rich morphology.
Our approach combines (1) a 3-level language-specific error taxonomy with 18 L2 categories capturing Kazakh-specific
phenomena such as vowel harmony violations and affix chain errors, (2) a dual-model architecture pairing
NLLB-200-distilled-600M seq2seq correction with an XLM-RoBERTa edit tagger using data-derived tags, and
(3) a multi-reference evaluation benchmark with 5 correction strategies.
We construct a mixed training dataset of synthetic (GPT-4o generated) and organic (Wikipedia edits, social media)
error-correction pairs, totaling approximately 87K examples.
Experiments show that the dual-model cascade outperforms either component alone, and morpheme-aware input
representation provides consistent gains on morphosyntactic error categories.
\end{abstract}

\section{Introduction}
\label{sec:intro}

% Motivation: GEC for agglutinative languages is understudied
% Kazakh: 15M+ speakers, Turkic family, rich morphology
% Challenges: affix chains, vowel harmony, limited annotated data
% Contributions: taxonomy, dual-model, multi-ref benchmark, first Kazakh GEC

\section{Related Work}
\label{sec:related}

\subsection{GEC for Morphologically Rich Languages}
% GECTurk (Turkish), KoGEC (Korean), Arabic GEC
% SCRIPT: seq2seq + edit models

\subsection{Multilingual Seq2Seq for GEC}
% NLLB-200, mBART as GEC backbone
% "Monolingual translation" approach

\subsection{Edit-Based GEC}
% GECToR, PIE, LaserTagger
% Data-derived vs hand-crafted tags

\section{Kazakh Error Taxonomy}
\label{sec:taxonomy}

% 3-level hierarchy: L1 (3) / L2 (18) / L3 (~40)
% L1: Orthography, Morphosyntax, Syntax \& Discourse
% Kazakh-specific: vowel harmony, affix chains, agglutinative morphology
% Table: full taxonomy with examples

\section{Data}
\label{sec:data}

\subsection{Synthetic Data Generation}
% GPT-4o taxonomy-tagged corruption
% 18 L2-balanced generation

\subsection{Organic Data Sources}
% Wikipedia kk edit history extraction
% Social media GPT-4o correction

\subsection{Data Mixing and Statistics}
% Deduplication, identity examples
% Train/val/test splits
% Distribution table

\section{Method}
\label{sec:method}

\subsection{NLLB-200 Seq2Seq Baseline}
% "Monolingual translation" kaz\_Cyrl $\to$ kaz\_Cyrl
% Fine-tuning details

\subsection{Edit Tagger}
% XLM-RoBERTa token classifier
% Data-derived tags: \$KEEP, \$DELETE, \$REPLACE\_X, \$APPEND\_X
% Top-K=2000 vocabulary

\subsection{Morpheme Segmentation}
% Apertium-kaz wrapper
% Qwen-distilled char-level segmenter
% Pipe-delimited input: бала|лар|ға

\subsection{Dual-Model Inference Pipeline}
% Cascade: Tagger (high-confidence) $\to$ NLLB seq2seq
% Threshold selection

\section{Multi-Reference Benchmark}
\label{sec:benchmark}

% 700 sentences $\times$ 5 strategies (minimal, conservative, moderate, fluent, alternative)
% 2--4 unique references per sentence
% Evaluation: multi-ref Word F$_{0.5}$, GLEU, CER

\section{Experiments}
\label{sec:experiments}

\subsection{Experimental Setup}
% Hardware, hyperparameters, training details

\subsection{Main Results}

% \begin{table}[t]
% \centering
% \begin{tabular}{lcccc}
% \toprule
% Model & F$_{0.5}$ & GLEU & CER & EM \\
% \midrule
% NLLB baseline (synthetic) & & & & \\
% NLLB (full data) & & & & \\
% Edit tagger only & & & & \\
% Cascade pipeline & & & & \\
% + morpheme input & & & & \\
% \bottomrule
% \end{tabular}
% \caption{Main results on the multi-reference test set.}
% \label{tab:main}
% \end{table}

\section{Ablation Studies}
\label{sec:ablation}

% Morpheme vs no-morpheme
% Tagger threshold sweep
% Data mixture ablation
% Model size comparison

\section{Analysis}
\label{sec:analysis}

\subsection{Per-Category Performance}
% Per-L1 and per-L2 F$_{0.5}$ breakdown
% Where does the system struggle?

\subsection{Error Analysis}
% Common failure modes
% Qualitative examples

\section{Conclusion}
\label{sec:conclusion}

% Summary of contributions
% Limitations
% Future work: larger models, more languages, human evaluation

\bibliographystyle{acl_natbib}
\bibliography{references}

\appendix

\section{Full Error Taxonomy}
\label{app:taxonomy}

\section{Data Generation Prompts}
\label{app:prompts}

\section{Additional Results}
\label{app:results}

\end{document}
